\documentclass{article}
\usepackage[UTF8]{ctex}
\usepackage{amsmath,amssymb,amsfonts}
\usepackage{enumitem}
\usepackage{geometry}
\geometry{a4paper, margin=1in}
% === 错题本排版优化（不参与题目标识） ===
\setlength{\parindent}{0pt}
\setlength{\parskip}{0.35em plus 0.10em minus 0.08em}
\setlength{\emergencystretch}{3em}
\setlist{itemsep=0.25em, topsep=0.40em, parsep=0pt}
\allowdisplaybreaks[2]
% === 错题本排版优化结束 ===
\begin{document}
\section*{合肥工业大学2025年数学分析真题}
\section*{108．合肥工业大学2025年数学分析试卷}

\subsection{题目1}
1．（10 分）求极限 \(\lim _{n \rightarrow \infty}\left(\frac{1}{n}-\frac{2}{n}+\frac{3}{n}+\cdots+\frac{(-1)^{n-1} n}{n}\right)\) ．

\subsection{题目2}
2．（10 分）求极限 \(\lim _{x \rightarrow+\infty} e^{-x}\left(1+\frac{1}{x}\right)^{x^{2}}\) ．

\subsection{题目3}
3．（15 分）设 \(f(x)\) 在 \([2,+\infty)\) 上连续，且 \(\lim _{n \rightarrow \infty} f(\sqrt{n})=0\) ．证明： \(\lim _{x \rightarrow+\infty} f(x)\) 存在当且仅当 \(f(x)\) 在\([2,+\infty)\) 上一致连续．

\subsection{题目4}
4．（15 分）设 \(f(x)\) 在 \([0,1]\) 上可导，且 \(f(0)=0, f(1)=1\) 。证明：在 \((0,1)\) 内存在两个不同的点\(x_{1}, x_{2}\) ，使得

\[
\frac{1}{f^{\prime}\left(x_{1}\right)}+\frac{2024}{f^{\prime}\left(x_{2}\right)}=2025 .
\]

\subsection{题目5}
5．（15 分）设函数 \(f(x)\) 在 \([0,2]\) 上连续可导，且 \(f(0)=f(2)=1,\left|f^{\prime}(x)\right| \leq 1\) ，证明： \(\int_{0}^{2} f(x) \mathrm{d} x>1\) ．

\subsection{题目6}
6．（10 分）讨论级数 \(\sum_{n=1}^{\infty} \frac{(-1)^{n-1}}{n} \frac{a}{1+a^{n}}(a>0)\) 的绝对收敛和条件收敛性．

\subsection{题目7}
7．（15 分）设 \(\left\{f_{n}(x)\right\}\) 为 \([a, b]\) 上的 Riemann 可积函数列，且函数列 \(\left\{f_{n}(x)\right\}\) 在 \([a, b]\) 上一致收敛于\(f(x)\) ．证明：\(f(x)\) 在 \([a, b]\) 上 Riemann 可积，且 \(\lim _{n \rightarrow \infty} \int_{a}^{b} f_{n}(x) \mathrm{d} x=\int_{a}^{b} f(x) \mathrm{d} x\) ．

\subsection{题目8}
8．（15 分）抛物面 \(x^{2}+y^{2}=z\) 被平面 \(x+y+z=1\) 截成一个椭圆，求这个椭圆到原点的最长与最短距离。

\subsection{题目9}
9．（15 分）计算曲面积分

\[
\iint_{S} y z \mathrm{~d} y \mathrm{~d} z+\left(x^{2}+z^{2}\right) y \mathrm{~d} z \mathrm{~d} x+x y \mathrm{~d} x \mathrm{~d} y .
\]

其中 \(S\) 为曲面 \(y=4-\left(x^{2}+y^{2}\right)(y>0)\) ，在 \(x o z\) 面右侧部分外侧。

\subsection{题目10}
10．（15 分）设 \(z=f(x-y, x+y)+g(x+k y), f\) 具有二阶连续偏导数，\(g\) 具有二阶连续导数且 \(g^{\prime \prime}(x)\)不恒为0。如果

\[
\frac{\partial^{2} z}{\partial x^{2}}+2 \frac{\partial^{2} z}{\partial x \partial y}+\frac{\partial^{2} z}{\partial y^{2}}=4 f_{22}^{\prime \prime}
\]

求常数 \(k\) 的值．

\subsection{题目11}
11．（15 分）求函数 \(I(\alpha)=\int_{1}^{+\infty} \frac{\sin x \cos x}{x^{\alpha}} \mathrm{d} x\) 的定义域，并研究其在定义域上的连续性。
\end{document}
